\section{Race and Death}

Since we now have 50 people in the entire world interested in reading more of Evola, I am releasing another section. It may be of value to compare it to Mircea Eliade's description of the Legionary movement\footnote{\url{http://ia600700.us.archive.org/1/items/Liberty\_7/Liberty.pdf}}, from which the following quote is taken.

\begin{quotex}
Being a profoundly Christian movement, justifying its doctrine on the spiritual level above all -- legionarism encourages and is built upon liberty. You adhere to legionarism because you are free, because you decided to overcome the iron circles of biological determinism (fear of death, suffering, etc.) and economic determinism (fear of becoming homeless). The first gesture a legionnaire will display is one born out of total liberty: he dares to free himself of spiritual, biological and economic enslavement. No exterior determinism can influence him. The moment he decides to be free is the moment all fears and inferiority complexes instantaneously disappear. He who enters the Legion forever dons the shirt of death. That means that the legionnaire feels so free that death itself no longer frightens him. If the Legionnaire nurtures the spirit of sacrifice with such passion, and if he has proven to be capable of making sacrifices these bear witness to the unlimited liberty a legionnaire has gained. ``He who knows how to die will never be a slave.'' And this doesn't concern only ethnic or political enslavement -- but firstly, spiritual enslavement. If you are ready to die, no fear, weakness, shyness can enslave you. Making peace with death is the most total liberty man can receive on this Earth.
\end{quotex}


The Evola text now follows. Note again that he opposes biological determinism and insists that the personality transcends all elements of race.

\begin{center}* * *\end{center}


Given how much has now been clarified, we will now pause a moment to make precise, in particular, the limits of the personality's membership in a race. Let us immediately say what is, in this respect, the view that is unacceptable from the traditional point of view. It is that of those who, once race is conceived as a purely biological-human, historical, and only terrestrial entity, holds that in such an entity lies the purpose of every being that belongs to it, that nothing exists higher than race, since race is the wellspring of every value, and the idea of a task and a super-terrestrial destination of the individual is illusory and pernicious: ``to remain faithful to the earth and to the race.''

We have already dealt with this conception and criticized it more than once. In the face of it, after all, one can resort to the racialist criterion for the appreciation of the ``truth''; depending on the various ``races of the spirit'' we have, particularly, as many conceptions of the same race, and there is no doubt that only for a telluric race can that conception be ``true'', only the telluric man can suppose such limited horizons as absolute. In this telluric vision of the race, the mentioned supposition of those neopagan racialists also comes back, according to whom the only conceivable immortality would be that of survival in blood, in earthly descendants.

It is true that similar positions, today, are presented less according to a theoretic value than according to a practical and political one. With it we aim to consolidate the unity of the people-race and to concentrate every spiritual energy of the individual member in the temporal and historical tasks that this being has to resolve. But it is otherwise true that the ancient Aryan civilizations, even in regard to earthly, heroic and political realizations, had their greatness without, however, perceiving the need of resorting to these myths, recognizing instead rather different truths. It is clearly evident, in fact, that our view of race carries back to the pitryana, to the ``way of the South'' and that it is opposed to the ``divine way of the North'' (devayana) that stood alone to define the highest Aryan ideal.

The theory of ``double heredity'' reconnects to such ideals. Personality, we said, is not exhausted in historical-biological heredity or horizontal heredity. It appears rather as a principle that, while manifesting itself in the race (here, as always, race in the restricted sense), in itself stands beyond race, therefore cannot be exhausted in it. To recognize race, as by now has been clarified in principle, does not mean to damage the personality: the personality owes the living and articulated material for its specific expression, for its self-manifesting and acting, to the race and as much as earthly heredity collects. In that there is a conditionality that however is not passive and one-sided. Every individual also reacts on race and on his heredity on the basis of his own most intimate nature, processes the substance in which he manifests himself, ultimately the form, and it is such that interracial differentiation and that diverse purity or completeness of types is realized, which we previously mentioned and now return to, in regards to its social reflexes. It is a giving rather than a receiving. In the points in which a supreme balance and a supreme adequacy are reached (a balance, according to our tripartite view of the various components of the true race), we have as a peak, beyond which the personality has nowhere to go --- it has nowhere to go on the horizontal, terrestrial line. To this line its work, its being, and, physiologically, its lineage remains and belongs. But the personality itself, if it has reached such a peak, is ``free'' and can now turn to a properly supernatural perfection.

This is exactly the most ancient Aryan conception, relative to which it does not belong specifically to the group of spiritual leaders, a conception that is also found again in the views and various legends of the Western Middle Ages. The dharma is established, the scrupulous observance of worldly laws of race, caste, and so on, up to a full adaptation. Such laws require also the assurance of a lineage: life, which is received at birth, before death must be returned with its own imprint to another being, and it is for this reason that the first-born son is called ``the son of duty''. After that, after the ``active life'', according to the Aryan law, he must withdraw to an ascetic-contemplative life. And the Ario-Iranian saying is rather expressive that recalls the true task to be, not just to procreate on the horizontal plane of terrestrial descent, but also towards the higher, on the vertical ascending direction. In the Western religion all these views were confused. Above all, what is of pertinence to the active life was separated violently from that which is the contemplative life and almost always the truly traditional solutions were forgotten, according to which the law, that in not of this earth, prolongs, completes and empowers that of this earth. But still more harmful than such confusions are the little noted telluric racist views, in case they have been taken seriously and have a future. According to the traditional teaching of the Aryan peoples, it remains instead firm that the supernatural is essentially the goal and the dignity of the personality: this goal, therefore, acts as the highest driving force and the deepest animating force in the bosom to the expression that race gives to the personality, thereby uplifting simultaneously the race, up to a limit beyond which, after having left a sign of greatness, the same force liberates itself and tends to make it so that death is precisely a finalization, a \emph{telos}, a new birth, the third birth of the indo-Aryan teaching.

Only the mediocre and the losers, among the beings who do not know how to realize the law and earthly duty down to their foundation, can think that they don't have an afterlife, that they have for a destiny a redissolving in the mixed vitality of the race, in the collective and earthly substance of blood and heredity, only in that way surviving (in a somewhat relative meaning of the word) the destruction of their physical individuality, and transmitting to others the task for which they were inadequate.


\flrightit{Posted on 2011-09-06 by Aeneas}

\begin{center}* * *\end{center}

\begin{footnotesize}
\begin{sffamily}
\texttt{HOO on 2011-09-06 at 22:20 said:}

This translation work is appreciated.

Although this was already our doctrine, the text, e.g. because of the way it is written, enlightens our own understanding. In the superior or complete worldview there is no divorce between holistic racialism and spirituality.

Meta-physical, extraphysical, or supernatural racism is the expression of the core of the religion of the Hyperborean-derived peoples --- liberty, hierarchy, personality; differentiation \& unity; --- and vice versa.

\hfill

\texttt{Jason on 2011-09-07 at 14:57 said:}

Cologero, in regards to your disappointment that so few have expressed interest in reading Evola, Traditionalism as such is meant for an elite, an elite is never a large group, and in this modern world the number of elite souls is much smaller than it would ever normally have been... ... ... .we do not know what the future will bring, and perhaps in the future there will be many who will be thankful to you for bringing out this important book in English. It is important that THE DOCTRINE OF RACE be translated if only to debunk the various ``Neo-Nazi'' interpretations of Evola common in many circles and to show the Baron's thoughts as they truly were. I for one, an admirer of Evola who cannot read Italiano, am very grateful for this project.


\hfill

\end{sffamily}
\end{footnotesize}

